% !PRINTED 619
\weberpage{619}
%<<CONT
the figures of $c$ unchanged, while it is nevertheless not the identical substitution, because $c_1^h$ does not leave all figures unchanged. But this is not possible when $h>1$. Hence either $\pi$ must leave one figure unchanged, or it must consist of cycles of equal degree. But since $n$ is a prime, $\pi$ can in this case form only a single cycle of degree $n$.

Thus $P$ contains, besides the identical permutation, only cyclic permutations of degree $n$, which we shall denote by $\gamma$, and permutations $\varkappa$ which leave one figure unchanged and moreover consist of cycles of equal degree.

Each of the permutations $\varkappa$ by which the figure 0 is left unchanged we shall denote by $\varkappa_0$. Likewise $\varkappa_1, \varkappa_2 \ldots \varkappa_{n-1}$ denote the permutations $\varkappa$ which leave the figures 1, 2 \ldots $n-1$ unchanged. Since $P$ is transitive, there must be just as many $\varkappa_0$ as $\varkappa_1$, as $\varkappa_2$, etc.

For if $\pi$ is a permutation which carries 0 into 1, then every $\pi^{-1}\varkappa_0\pi$ is a $\varkappa_1$, and conversely every $\pi\varkappa_1\pi^{-1}$ is a $\varkappa_0$, and likewise for the remaining figures. If then $\mu$ is the number of the $\varkappa_0$, $\nu$ the number of the $\gamma$, $m$ the degree of $P$, then, since the identical permutation is also to be added, we have

\weq{11}{m = \mu n + \nu + 1.}

Now the $\varkappa_0$, together with the identical permutation, form a group $Q$ of degree $\mu+1$ and a divisor of $P$, namely the totality of all permutations of $P$ which leave 0 unchanged. If then $\pi_1, \pi_2 \ldots \pi_{n-1}$ are permutations in $P$ which carry 0 into $1, 2 \ldots n-1$, then $Q\pi_1$ is the totality of all those permutations of $P$ which carry 0 into 1, and so forth.

We thus obtain the decomposition of $P$ into the subsidiary groups

\[
P = Q + Q\pi_1 + Q\pi_2 + \cdots + Q\pi_{n-1},
\]

from which follows

\weq{12}{m = n(\mu+1),}

and hence from (11) and (12) $\nu = n-1$.

There are therefore $n-1$ and no more cyclic permutations of degree $n$ in $P$, and these consequently form, again together with the unit, a cyclic group

\[
C = 1, \gamma, \gamma^2 \ldots \gamma^{n-1},
\]

since with $\gamma$ all powers of $\gamma$ likewise occur in $P$.

Now if $\gamma$ is a cyclic permutation of degree $n$, then every $\pi^{-1}\gamma^h\pi$ is also cyclic, and hence, if $\pi$ belongs to $P$
%CONT>>
